\label{sec:lit}

\subsection{Windows Event Log Anomaly Detection}

Windows Event Logs contain semantically rich, structured records of system activity.
He~\textit{et al.}\ \cite{he2017drain} introduced Drain, an online log-parsing tree that
maps raw log messages to templates in $O(\text{log}~n)$ time; their approach is directly
applicable to tokenising Windows Event Log descriptions at scale. Du~\textit{et al.}\
\cite{du2017deeplog} modelled log key sequences as a workflow using LSTM networks, detecting
anomalies as deviations from learned sequential patterns. Their CCS~2017 evaluation
demonstrated low false-positive rates on the HDFS log dataset. Zhao~\textit{et al.}\
\cite{zhao2019loganomaly} extended this to LogAnomaly, capturing both sequential and
quantitative anomalies via template-based semantic vectors. Garcia-Teodoro~\textit{et al.}\
\cite{garcia2009anomaly} provided a foundational survey of anomaly-based network intrusion
detection, whose classification taxonomy—statistical, knowledge-based, and
machine-learning—applies equally to host-based log analysis. Collectively, these works
demonstrate the viability of sequential log modelling but leave the ICS-specific
event-identifier semantics unexploited.

\subsection{ICS/SCADA Intrusion Detection}

Industrial control system security demands domain-adapted detection strategies. Carcano
\textit{et al.}\ \cite{carcano2011scada} proposed a multidimensional critical-state
analysis that flags SCADA command sequences deviating from learned safe states, achieving
100\% detection on a simulated water treatment testbed. Goldenberg and Wool
\cite{goldenberg2013modbus} built deterministic finite automata from Modbus/TCP network
traffic, detecting protocol violations with zero false positives in a power utility
environment. Mitchell and Chen \cite{mitchell2014survey} surveyed IDS techniques for
cyber-physical systems, categorising approaches by detection method (misuse vs.\ anomaly)
and system layer (network, process, device). Goh~\textit{et al.}\ \cite{goh2016swat}
released the Secure Water Treatment (SWaT) dataset, enabling reproducible evaluation of
process-level anomaly detectors. Zhu~\textit{et al.}\ \cite{zhu2011scada} formalised a
taxonomy of cyber attacks on SCADA systems across all Purdue model levels, which informs
the attack patterns modelled in this work. A key gap identified across this body of work is
the absence of Windows host-log analysis as a detection layer within the ICS hierarchy.

\subsection{AI/ML Methods for Intrusion Detection}

Machine learning for intrusion detection has a long history. Liao~\textit{et al.}\
\cite{liao2013ids} reviewed supervised, unsupervised, and hybrid methods, noting that
ensemble methods consistently outperform single classifiers on tabular security data.
Breiman's Random Forest \cite{breiman2001}—a bagged ensemble of decision trees—remains a
strong baseline due to its resistance to overfitting, built-in feature importance ranking,
and native support for class imbalance through per-class sample weighting. Chen and
Guestrin's XGBoost \cite{chen2016xgboost} extended gradient boosting to large-scale tabular
problems, achieving state-of-the-art results on the KDD~Cup~1999 network intrusion
dataset. Tavallaee~\textit{et al.}\ \cite{tavallaee2009kdd} revealed severe redundancy
in the KDD~Cup~1999 benchmark, motivating the community towards domain-specific synthetic
datasets. Sommer and Paxson \cite{sommer2010outside} argued that the closed-world
assumption of supervised ML makes deployment in open operational environments inherently
difficult, underscoring the need for realistic synthetic training data and rigorous
evaluation protocols. NIST SP~800-82 \cite{stouffer2015nist} provides the governing
reference for ICS security controls, including detection and response requirements that
bound the target performance metrics for this work.
